% GENERATED FILE. Edit canonical lessons, then run npm run generate:books.
% canonical-source-hash: fnv1a64:1cc87b0ac173830a
% canonical-lessons: ZH-C150-guoqu, ZH-W150-guo-pass, ZH-C150-guo-ever, ZH-C150-houlai, ZH-C150-ranhou, ZH-C150-conglai, ZH-C150-practice

\chapter{Have You Ever Been?}
\label{ch:zh-have-you-ever-been}
\hlchaptermodality{150}

\begin{chapteropening}
\textbf{By the end of this chapter:} \emph{I can ask whether someone has ever done something, and answer.}

Asking whether someone has ever done something with \zh{过} after the verb, and telling what came next with \zh{后来} and ránhòu, built on one new written character: \zh{过}.
\end{chapteropening}

\section[\texorpdfstring{\zh{过去} (guòqù)}{过去 (guòqù)}]{\zh{过去} (guòqù) — in the past; to go over}

\label{lesson:ZH-C150-guoqu}



\begin{quote}
A new chapter, and it opens on time that has already gone by.

\begin{itemize}
\raggedright
  \item \emph{Recall:} say \zh{我正在吃饭} — \emph{wǒ zhèngzài chī fàn} — \textbf{R1}, the chapter you have finished
  \item \emph{Recall:} say \zh{一边} — \emph{yìbiān}, at the same time — \textbf{R1}
  \item \emph{Recall:} say \zh{以前} — \emph{yǐqián}, before — \textbf{R2}, a few chapters back
\end{itemize}
\end{quote}

\subsection*{You'll want to know: \zh{过去}}
\textbf{\zh{过去}} — \emph{guòqù} — \textquotedblleft{}in the past; to go over (there)\textquotedblright{}.

\begin{quote}
\textbf{\zh{过去} \zh{我家在大学里边}} \emph{guòqù, wǒ jiā zài dàxué lǐbian} — In the past, our home was on the university grounds.
\end{quote}

\begin{quote}
\textbf{\zh{我过去}} \emph{wǒ guòqù} — I'll go over there.
\end{quote}

The next lesson puts the first character in your hand. It is the guò of \emph{guònián} and \emph{nánguò}.

\begin{cousinweb}
\textbf{\zh{过去}} (\emph{guòqù}) is built from \zh{过} 'to pass' and \zh{去} 'to go'. Both are inherited from Middle Chinese, \emph{kwaH} and \emph{khjoH} (Guangyun, in Baxter's transcription), glossed 'to go beyond; to exceed' and 'to leave'. 'The past', time that has gone by, is the pair's own sense.
\end{cousinweb}

\begin{grammarlens}[title={What you've built}]
\zh{过去} and \zh{以前} both look back. \zh{以前} is before some point; \zh{过去} is the past at large, set against now, \emph{xiànzài}.

As a verb, \zh{过去} is going over there, away from the speaker.
\end{grammarlens}

\subsection*{Your turn}
\begin{itemize}
\raggedright
  \item \emph{Say it:} \emph{guòqù}
  \item \emph{Say it:} \emph{guòqù}, once more
  \item \emph{Say it:} \zh{过去} \zh{我家在大学里边} — \emph{guòqù, wǒ jiā zài dàxué lǐbian} — In the past, our home was on the university grounds.
  \item \emph{Say it:} \zh{我过去} — \emph{wǒ guòqù} — I'll go over there.
  \item \emph{Recall:} say \zh{以前}, then \zh{过去}
\end{itemize}

\begin{tcolorbox}[breakable,colback=teal!4,colframe=teal!35!black,title={Before you move on}]
What does \zh{过去} mean? (\textquotedblleft{}In the past; to go over\textquotedblright{}.) And now? (\emph{xiànzài}.)
\end{tcolorbox}
\section[\texorpdfstring{\zh{过} (guò)}{过 (guò)}]{\zh{过} — \zh{寸}, then the walking piece last}

\label{lesson:ZH-W150-guo-pass}



\begin{quote}
The word first, then the character inside it.

\begin{itemize}
\raggedright
  \item \emph{Recall:} say \emph{guòqù} — \textbf{R1}, one lesson back, before it decays
\end{itemize}

\begin{quote}
\zh{对} \zh{时} \zh{这} \zh{进}
\end{quote}

\emph{Read it:} each one aloud, with nothing said first — \emph{duì}, \emph{shí}, \emph{zhè}, \emph{jìn}
\end{quote}

\subsection*{Script — \zh{寸}, then the walking piece last}
\hlblockfigure{\detokenize{figures/ZH-W150-guo-pass-filmstrip.pdf}}{How \zh{过} is written, stroke by stroke}

\begin{quote}
\zh{过}
\end{quote}

\textbf{Two pieces you own, in the order you know from \zh{这}.}

\textbf{\zh{寸}, three strokes}, on the upper right: the bar, the vertical that hooks to the left at its foot, then the dot. It is the right half of \zh{对} and of \zh{时}.

\textbf{\zh{辶}, three strokes}, last: the dot at the upper left; the short stroke that turns down, turns again and falls to the left in one movement; then the long base, sweeping out to the right under \zh{寸}.

3 + 3 = six. The walking piece goes in last, as in \zh{这} and \zh{进}, though it sits under everything.

\subsection*{Your turn}
\begin{itemize}
\raggedright
  \item \emph{Write it:} \zh{过}, all six, with the model in front of you
  \item \emph{Cover:} write it again from memory, and count the pieces, not the strokes
\end{itemize}

\begin{tcolorbox}[breakable,colback=teal!4,colframe=teal!35!black,title={Before you move on}]
How many strokes in \zh{过}? (\textbf{Six}.) Which piece first? (\textbf{\zh{寸}}.) Which character you write puts the walking piece last in the same way? (\textbf{\zh{这}, and \zh{进}}.)

Source: \href{https://raw.githubusercontent.com/chanind/hanzi-writer-data/68d10a4b21150cae5e1ebbd223eed289cf32d90c/data/\%E8\%BF\%87.json}{Hanzi Writer Data \zh{过}, pinned stroke paths and medians}.
\end{tcolorbox}
\section[\texorpdfstring{\zh{你去过中国吗} (nǐ qù guo Zhōngguó ma)}{你去过中国吗 (nǐ qù guo Zhōngguó ma)}]{\zh{你去过中国吗} — \zh{过} for what you have done}

\label{lesson:ZH-C150-guo-ever}



\begin{quote}
The character you have written turns up again, light, straight after a verb.

\begin{itemize}
\raggedright
  \item \emph{Read it:} \zh{过去}, aloud — \emph{guòqù}
  \item \emph{Recall:} say \zh{我去了中国} — \emph{wǒ qù le Zhōngguó}, I went to China
  \item \emph{Recall:} say \zh{我没去} — \emph{wǒ méi qù}, I didn't go
\end{itemize}
\end{quote}

\begin{grammarlens}[title={Grammar Lens: \zh{了} says it happened; \zh{过} says it has happened at least once}]
\begin{quote}
\textbf{\zh{过} straight after a verb says this has happened at least once, at some time. Here it is said light: \emph{guo}.}
\end{quote}

\begin{quote}
\textbf{\zh{你去过中国吗}} \emph{nǐ qù guo Zhōngguó ma} — Have you ever been to China? \textbf{\zh{我去过}} \emph{wǒ qù guo} — I have.
\end{quote}

The negative keeps \zh{过} and puts \zh{没} in front; \zh{不} has no place here, because the question is about what has happened:

\begin{quote}
\textbf{\zh{我没去过}} \emph{wǒ méi qù guo} — I've not been.
\end{quote}

Set it beside \zh{了}. \zh{我去了中国} reports one trip; \zh{我去过中国} says China is somewhere you have been, at some point in your life:

\begin{quote}
\textbf{\zh{我去了中国}} \emph{wǒ qù le Zhōngguó} — I went to China. \textbf{\zh{我去过中国}} \emph{wǒ qù guo Zhōngguó} — I've been to China.
\end{quote}
\end{grammarlens}

\subsection*{Your turn}
Say these aloud:

\begin{itemize}
\raggedright
  \item \zh{你去过中国吗} — \emph{nǐ qù guo Zhōngguó ma} — Have you ever been to China?
  \item \zh{我去过} — \emph{wǒ qù guo} — I have.
  \item \zh{我没去过} — \emph{wǒ méi qù guo} — I've not been.
  \item \zh{你看过这本书吗} — \emph{nǐ kàn guo zhè běn shū ma} — Have you read this book?
\end{itemize}

\begin{tcolorbox}[breakable,colback=teal!4,colframe=teal!35!black,title={Before you move on}]
What does \zh{过} after a verb say? (\textbf{It has happened at least once}.) Which negative goes with it? (\textbf{\zh{没}: \zh{我没去过}}.) Which says one trip that happened, \zh{了} or \zh{过}? (\textbf{\zh{了}}.)
\end{tcolorbox}
\section[\texorpdfstring{\zh{后来} (hòulái)}{后来 (hòulái)}]{\zh{后来} (hòulái) — afterwards; later}

\label{lesson:ZH-C150-houlai}



\begin{quote}
Before the new one: read the last lesson's question, then the day you wrote \zh{后} for.

\begin{itemize}
\raggedright
  \item \emph{Read it:} \zh{你去过中国吗}, aloud — \emph{nǐ qù guo Zhōngguó ma}
  \item \emph{Recall:} say \zh{后天} — \emph{hòutiān}, the day after tomorrow — \textbf{R2}, a few chapters back
\end{itemize}
\end{quote}

\subsection*{You'll want to know: \zh{后来}}
\textbf{\zh{后来}} — \emph{hòulái} — \textquotedblleft{}afterwards; later (in the past)\textquotedblright{}.

\begin{quote}
\textbf{\zh{后来我们去了大学}} \emph{hòulái wǒmen qù le dàxué} — Afterwards we went to the university.
\end{quote}

\begin{cousinweb}
\textbf{\zh{后来}} (\emph{hòulái}) is built from \zh{后} 'after' and \zh{来} 'to come', both inherited from Middle Chinese, \emph{huwX} and \emph{loj} (Guangyun, in Baxter's transcription): what came after.
\end{cousinweb}

\begin{grammarlens}[title={What you've built}]
\zh{后来} looks back: it follows something that already happened and says what came next. For what will come later, Mandarin says \emph{yǐhòu}: \emph{yí ge xiǎoshí yǐhòu}.

One more way to set a time in the past: \emph{de shíhou} after a clause turns it into a when, the shíhou of \emph{xiǎoshíhou}:

\begin{quote}
\emph{wǒ xiǎo de shíhou, wǒ jiā zài Zhōngguó, hòulái wǒmen lái zhèr le} — When I was small my family lived in China; later we came here.
\end{quote}
\end{grammarlens}

\subsection*{Your turn}
\begin{itemize}
\raggedright
  \item \emph{Say it:} \emph{hòulái}
  \item \emph{Say it:} \emph{hòulái}, once more
  \item \emph{Say it:} \zh{后来我们去了大学} — \emph{hòulái wǒmen qù le dàxué} — Afterwards we went to the university.
  \item \emph{Say it:} \emph{wǒ xiǎo de shíhou, wǒ jiā zài Zhōngguó} — When I was small my family lived in China.
  \item \emph{Recall:} say \zh{后天}, then \zh{后来}, and hear the same hòu
\end{itemize}

\begin{tcolorbox}[breakable,colback=teal!4,colframe=teal!35!black,title={Before you move on}]
What does \zh{后来} mean? (\textquotedblleft{}Afterwards; later, in the past\textquotedblright{}.) And for later still to come? (\emph{yǐhòu}.)
\end{tcolorbox}
\section[ránhòu]{ránhòu — then; after that}

\label{lesson:ZH-C150-ranhou}



\begin{quote}
Before the new one: last lesson's word, then a word that ends the same way.

\begin{itemize}
\raggedright
  \item \emph{Recall:} say \zh{后来} — \emph{hòulái}, afterwards — \textbf{R1}, one lesson back
  \item \emph{Recall:} say \emph{yǐhòu}, after — \textbf{R3}, durable
\end{itemize}
\end{quote}

\subsection*{You'll want to know: ránhòu}
\emph{ránhòu} — \textquotedblleft{}then; after that\textquotedblright{}.

\begin{quote}
\emph{wǒ xiān chī fàn, ránhòu kàn shū} — First I eat, then I read.
\end{quote}

Its first character comes later; for now it is a word for the ear.

\begin{cousinweb}
\emph{ránhòu} is built from rán 'so; thus' and \zh{后} 'after'. Both are inherited from Middle Chinese, \emph{nyen} and \emph{huwX} (Guangyun, in Baxter's transcription); the Guangyun glosses \emph{nyen} as a particle, 'thus', and 'then' for the pair is assumed here, not quoted.
\end{cousinweb}

\begin{grammarlens}[title={What you've built}]
\emph{ránhòu} joins two steps in order, the second after the first, now or in the past: \emph{wǒ xiān chī fàn, ránhòu kàn shū}.

\zh{后来} only looks back at what happened next; \emph{ránhòu} can plan tomorrow too: \emph{míngtiān wǒ qù dàxué, ránhòu qù kàn shū}.
\end{grammarlens}

\subsection*{Your turn}
\begin{itemize}
\raggedright
  \item \emph{Say it:} \emph{ránhòu}
  \item \emph{Say it:} \emph{ránhòu}, once more
  \item \emph{Say it:} \emph{wǒ xiān chī fàn, ránhòu kàn shū} — First I eat, then I read.
  \item \emph{Recall:} say \zh{后来}, then \emph{ránhòu}
\end{itemize}

\begin{tcolorbox}[breakable,colback=teal!4,colframe=teal!35!black,title={Before you move on}]
What does \emph{ránhòu} mean? (\textquotedblleft{}Then; after that\textquotedblright{}.) Which of the two can plan tomorrow, \zh{后来} or \emph{ránhòu}? (\emph{\textbf{ránhòu}}.)
\end{tcolorbox}
\section[cónglái]{cónglái — ever; never once}

\label{lesson:ZH-C150-conglai}



\begin{quote}
Before the new one: last lesson's word, then the answer to have you ever.

\begin{itemize}
\raggedright
  \item \emph{Recall:} say \emph{ránhòu}, then — \textbf{R1}, one lesson back
  \item \emph{Recall:} say \zh{我没去过} — \emph{wǒ méi qù guo}, I've not been — \textbf{R1}, this chapter
\end{itemize}
\end{quote}

\subsection*{You'll want to know: cónglái}
\emph{cónglái} — \textquotedblleft{}ever\textquotedblright{}; before \zh{没}, \textquotedblleft{}never once\textquotedblright{}.

\begin{quote}
\emph{wǒ cónglái méi qù guo} — I've never once been there.
\end{quote}

Its first character comes later; for now it is a word for the ear.

\begin{cousinweb}
\emph{cónglái} is built from cóng 'from' and \zh{来} 'to come': from then until now. Both are inherited from Middle Chinese, \emph{dzjowng} and \emph{loj} (Guangyun, in Baxter's transcription), \emph{dzjowng} glossed 'to go to; to follow'. 'Ever' is the pair's own sense, not the gloss's.
\end{cousinweb}

\begin{grammarlens}[title={What you've built}]
\emph{cónglái} stands in front of \zh{没} … \zh{过} and makes it stronger: \zh{我没去过}, I've not been; \emph{wǒ cónglái méi qù guo}, not once.

It needs \zh{没} after it to say never; on its own it is rare.
\end{grammarlens}

\subsection*{Your turn}
\begin{itemize}
\raggedright
  \item \emph{Say it:} \emph{cónglái}
  \item \emph{Say it:} \emph{cónglái}, once more
  \item \emph{Say it:} \emph{wǒ cónglái méi qù guo} — I've never once been there.
  \item \emph{Say it:} \zh{我没去过} — \emph{wǒ méi qù guo} — I've not been.
  \item \emph{Recall:} say \emph{ránhòu}, then \emph{cónglái}
\end{itemize}

\begin{tcolorbox}[breakable,colback=teal!4,colframe=teal!35!black,title={Before you move on}]
What does \emph{cónglái méi} … \zh{过} say? (\textbf{Never once}.) What does it strengthen? (\textbf{\zh{我没去过}}.)
\end{tcolorbox}
\section[Practice]{\zh{你去过吗} — \zh{我没去过}}

\label{lesson:ZH-C150-practice}



\begin{quote}
The chapter closes here, with two people comparing where they have been.

\begin{itemize}
\raggedright
  \item \emph{Recall:} say \emph{cónglái}, ever — \textbf{R1}, one lesson back
  \item \emph{Recall:} say \emph{ránhòu}, then — \textbf{R1}
  \item \emph{Recall:} say \zh{后来} — \emph{hòulái}, afterwards — \textbf{R2}
\end{itemize}
\end{quote}

\subsection*{The exchange}
\begin{quote}
\textbf{A:} \zh{你去过中国吗}
\end{quote}

\begin{quote}
\emph{nǐ qù guo Zhōngguó ma} — Have you ever been to China?
\end{quote}

\begin{quote}
\textbf{B:} \zh{去过}
\end{quote}

\begin{quote}
\emph{qù guo} — I have.
\end{quote}

\begin{quote}
\textbf{B:} \zh{过去我家在中国}
\end{quote}

\begin{quote}
\emph{guòqù wǒ jiā zài Zhōngguó} — My family used to live there.
\end{quote}

\begin{quote}
\textbf{A:} \zh{后来呢}
\end{quote}

\begin{quote}
\emph{hòulái ne} — And afterwards?
\end{quote}

\begin{quote}
\textbf{B:} \zh{后来我们来这儿了}
\end{quote}

\begin{quote}
\emph{hòulái wǒmen lái zhèr le} — Afterwards we came here.
\end{quote}

\begin{quote}
\textbf{B:} \zh{你呢} \zh{你去过吗}
\end{quote}

\begin{quote}
\emph{nǐ ne, nǐ qù guo ma} — And you? Have you been?
\end{quote}

\begin{quote}
\textbf{A:} \zh{我没去过}
\end{quote}

\begin{quote}
\emph{wǒ méi qù guo} — I've not been.
\end{quote}

\begin{quote}
\textbf{B:} \zh{你开过车吗}
\end{quote}

\begin{quote}
\emph{nǐ kāi guo chē ma} — Have you ever driven a car?
\end{quote}

\begin{quote}
\textbf{A:} \emph{wǒ cónglái méi kāi guo} — Never once.
\end{quote}

\begin{quote}
\textbf{B:} \emph{wǒ xiān qù Zhōngguó, ránhòu qù Běijīng} — I'll go to China first, then to Beijing.
\end{quote}

\zh{去过} answers \zh{你去过中国吗} on its own: the verb and \zh{过}, nothing else. The no is \zh{没去过}, with \zh{过} kept.

\begin{tcolorbox}[breakable,skin=enhanced,colback=orange!6,colframe=orange!45!black,boxrule=0.5pt,arc=1mm,left=6pt,right=6pt,top=4pt,bottom=4pt,fonttitle=\bfseries,title={Read this}]
\emph{Read it:} the exchange through once without stopping, then answer the question below
\end{tcolorbox}

\begin{grammarlens}[title={What you've built — ever, and what came next}]
\zh{过} after a verb asks about a whole life; \zh{后来} and \emph{ránhòu} put events in order:

\begin{itemize}
\raggedright
  \item \textbf{\zh{你去过中国吗}} \emph{nǐ qù guo Zhōngguó ma} — have you ever been to China?
  \item \textbf{\zh{我没去过}} \emph{wǒ méi qù guo} — I've not been.
  \item \textbf{\zh{后来我们来这儿了}} \emph{hòulái wǒmen lái zhèr le} — afterwards we came here.
  \item \emph{wǒ xiān chī fàn, ránhòu kàn shū} — first I eat, then I read.
\end{itemize}

\textbf{Listen.} \emph{Nǐ qù guo Běijīng ma?}

\textbf{Speak.} Say one place you have been, and one you have not.

\emph{Write it:} \zh{过} — the right half of \zh{对} first, the walking piece last
\end{grammarlens}

\subsection*{Your turn}
\begin{itemize}
\raggedright
  \item \emph{Say it:} \zh{你去过中国吗} — \emph{nǐ qù guo Zhōngguó ma} — Have you ever been to China?
  \item \emph{Say it:} \zh{我没去过} — \emph{wǒ méi qù guo} — I've not been.
  \item \emph{Say it:} the whole exchange, both parts
  \item \emph{Write it:} \zh{你去过中国吗}
\end{itemize}

\begin{tcolorbox}[breakable,colback=teal!4,colframe=teal!35!black,title={Before you move on}]
Ask if someone has ever been to China. (\textbf{\zh{你去过中国吗}}.) Say you have not. (\textbf{\zh{我没去过}}.) Say what happened afterwards, with \zh{后来}. (\textbf{\zh{后来我们来这儿了}}.)
\end{tcolorbox}
